\documentclass[11pt]{article}
\usepackage[a4paper,margin=27mm]{geometry}
\usepackage{amsmath,amssymb,amsthm,mathtools,microtype,booktabs,tabularx}
\usepackage[T1]{fontenc}
\usepackage{lmodern}
\usepackage{xcolor}
\definecolor{linkblue}{RGB}{25,65,105}
\usepackage[colorlinks=true,linkcolor=linkblue,citecolor=linkblue,urlcolor=linkblue]{hyperref}
\usepackage{attachfile2}
\hypersetup{pdftitle={The irrationality exponent of pi through a two-case determinant argument},pdfauthor={Ryan Matthew Casper},pdfsubject={A reorganized proof with a global-generation interpolation argument}}
\newtheorem{theorem}{Theorem}[section]
\newtheorem{proposition}[theorem]{Proposition}
\newtheorem{lemma}[theorem]{Lemma}
\newtheorem{corollary}[theorem]{Corollary}
\theoremstyle{definition}
\newtheorem{definition}[theorem]{Definition}
\theoremstyle{remark}
\newtheorem{remark}[theorem]{Remark}
\newcommand{\C}{\mathbb C}
\newcommand{\Q}{\mathbb Q}
\newcommand{\Z}{\mathbb Z}
\newcommand{\N}{\mathbb Z_{\ge0}}
\newcommand{\ii}{\mathrm i}
\newcommand{\ord}{\operatorname{ord}}
\newcommand{\lcm}{\operatorname{lcm}}
\newcommand{\wt}{\operatorname{wt}}
\newcommand{\Span}{\operatorname{span}}
\newcommand{\len}{\operatorname{length}}
\newcommand{\Pcal}{\mathcal P}
\newcommand{\Jcal}{\mathcal J}
\newcommand{\OO}{\mathcal O}
\numberwithin{equation}{section}
\setlength{\parskip}{0.35em}
\setlength{\emergencystretch}{2em}
\title{The irrationality exponent of \(\pi\)\\through a two-case determinant argument}
\author{Ryan Matthew Casper}
\date{October 11, 2026\\Revised November 26, 2026}

\begin{document}
\maketitle
\begin{abstract}
We present OpenAI's proof that \(\mu(\pi)=2\), with a different argument
for the interpolation step. Given sufficiently good rational
approximations, interpolation supplies a nonzero determinant and
denominator clearing bounds its modulus below. We translate the rows
to exact logarithmic periods and bound each term in the expansion.
If a term has many rows of low transverse weight, there are few
transverse indices available to them. Equal Taylor degrees within an
index give equal rows. The surviving degrees are therefore distinct,
and their sum grows quadratically with the number of rows sharing the
index. If a term has many high-weight rows, powers of the approximation
errors give the saving. With suitable parameters, the resulting upper
bound contradicts the arithmetic lower bound.
For interpolation itself we keep the original curve inequality and
ordinary blow-up. Globally generated sections then give polynomials
whose local order exceeds their source degree, and repeated reduction
by their weighted initial forms lifts jets of weight below \(T\) in
source degree at most \(\kappa T+B\) for fixed \(\kappa<1\) and \(B\). The Lean
companion checks this interpolation argument and its use in the
proof that \(\mu(\pi)=2\).
\end{abstract}

\begingroup
\small
\setlength{\parskip}{0pt}
\tableofcontents
\endgroup
\clearpage

\section{The claim and the organization of the proof}
For an irrational real number \(x\), set
\[
 \mu(x)=\sup\left\{\nu>0:
 0<|x-p/q|<q^{-\nu}\text{ for infinitely many coprime }p,q\in\Z,
 \ q\ge2\right\}.
\]
We prove the following form of the upper bound.
\begin{theorem}\label{thm:main}
For every real \(\nu>2\), there is an integer \(Q(\nu)\) such that
\[
 |\pi-p/q|>q^{-\nu}\qquad(p\in\Z,\ q\in\Z,\ q\ge Q(\nu)).
\]
Consequently \(\pi\) is irrational and \(\mu(\pi)=2\).
\end{theorem}

For each fixed \(\nu>2\), the strict inequality rules out
approximations of that quality beyond a cutoff \(Q(\nu)\).
The cutoff is noneffective. A uniform lower bound of the form \(c/q^2\)
would be a stronger statement.

One determinant carries the whole contradiction. Suppose exponent-\(\nu\) approximations exist
at unbounded denominators. A weighted interpolation theorem supplies a
nonzero determinant \(\Delta_H\) made from monomials in the original
polynomial space. Denominator clearing turns a scaled form of this
determinant into a nonzero Gaussian integer (an element of
\(\Z[\ii]\), hence of modulus at least one). Undoing the scaling leaves
an explicit lower bound for \(|\Delta_H|\). To estimate the same
determinant from above, we translate its rows to the exact periods
\(2\pi\ii j\) and expand. Each term of that row expansion has one of two
forms:
\begin{enumerate}
\item Many transverse indices are small. There are only a few such
indices, so many rows share the same index. Within one index, a repeated
Taylor degree gives equal coefficient rows, and the term vanishes. The
degrees that survive are therefore distinct, with a total at least
quadratic in the number of rows sharing the index. Taylor decay turns
that large total degree into an analytic saving.
\item Few transverse indices are small. The remaining rows have enough
transverse weight for powers of the approximation errors to give the
saving after we account for the logarithm tails.
\end{enumerate}
We choose the parameters so that both savings exceed the arithmetic
cost, with room for every error term and for the number of terms in the
expansion.

We take the number-theoretic half first.
Sections~\ref{sec:determinant}--\ref{sec:parameters} assume only the
interpolation theorem stated below and derive the contradiction from it.
Sections~\ref{sec:curve}--\ref{sec:bridge} then prove that theorem, so
the reader can see how the exponent bound follows before working through
the algebraic geometry. The later proof uses standard facts about
isolated-zero B\'ezout bounds, line bundles, ordinary blow-ups and
ampleness.

The curve inequality, ordinary blow-up, and determinant estimates come
from OpenAI~\cite{openai}. OpenAI obtains interpolation by pushing
powers of the ideal sheaf of the exceptional divisor down to the
original compactification and applying Serre vanishing. We use global
generation and finite polynomial reduction in Section~\ref{sec:bridge}
to account for the source degree of a jet lift factor by factor. Each
reduction factor contributes more local weight than it uses in source
degree, and the lifts of the low-weight remainders have a fixed degree
bound. The determinant comparison is in the tradition of interpolation
determinants \cite{laurent1991,laurent2000}.

\subsection{Notation and the interpolation input}
All logarithms are natural. We compare the space of polynomials of
bounded weighted degree with the truncated local expansions to be
realized at each center. For a
positive rational weight vector
\(W=(w_0,\ldots,w_m)\), let \(\Pcal_W(H)\) be the span over \(\C\)
of the monomials
\[
 Y^hX^\alpha=Y^h\prod_{i=1}^mX_i^{\alpha_i},
 \qquad w_0h+\sum_{i=1}^mw_i\alpha_i\le H.
\]
For positive local weights \(V=(v_0,\ldots,v_m)\), the local order
of a nonzero series \(F\in\C[[t,u_1,\ldots,u_m]]\) is
\[
 \ord_V F=\min\left\{v_0s+\sum_{i=1}^m v_i\beta_i:
                     [t^su^\beta]F\ne0\right\},\qquad
 \ord_V 0=+\infty.
\]
We write \(\ord_P\) for ordinary order on a branch \(P\). The jet
space \(\Jcal_V(H)\) is the quotient by the ideal of series of local
order at least \(H\). It keeps the finitely many coefficients of
weight less than \(H\). Multi-indices have nonnegative integer entries,
and inequalities between them are componentwise. We abbreviate
\(w\alpha=\sum_{i=1}^m w_i\alpha_i\).
Interpolation will ask whether source polynomials in \(\Pcal_W(H)\) can
realize prescribed elements of these truncated spaces at finitely many
centers at once.

\begin{theorem}[Separated-weight interpolation]\label{thm:interpolation}
Fix integers \(m,K\ge1\) and positive rational \(w_0,v_0,\theta\)
such that
\begin{equation}\label{eq:volume}
 0<\theta<1,\qquad K\theta^m<1,\qquad
 K(w_0/v_0)\theta^m<1.
\end{equation}
There are successive lower thresholds for positive rational weights
\(w_1,\ldots,w_m\), each depending only on these data and the preceding
weights, with the following property. Put
\(V=(v_0,w_1/\theta,\ldots,w_m/\theta)\). For any centers
\(a_j=(1,c_{j1},\ldots,c_{jm})\), \(0\le j<K\), whose coordinates
are pairwise distinct across \(j\) for each fixed \(i\), the map
\begin{equation}\label{eq:interpolation}
 \Pcal_W(H)\longrightarrow\bigoplus_{j=0}^{K-1}\Jcal_V(H),\qquad
 P\longmapsto
 \left(P(1+t,(c_{ji}+\log(1+t)+u_i)_{i=1}^m)\right)_j
\end{equation}
is surjective for all sufficiently large real \(H\).
The threshold for \(H\) may depend on the centers; the successive
weight thresholds do not.
\end{theorem}

Because the successive weight thresholds do not depend on the centers,
we can first choose good approximation denominators successively to meet
those thresholds. The centers are then fixed, so we can take \(H\) past
their interpolation threshold. The original formulation requires
sufficiently divisible integers \(H\). Section~\ref{sec:bridge} gives
interpolation for all sufficiently large real \(H\). Either form
suffices for the determinant argument.

\section{One determinant and two analytic alternatives}\label{sec:determinant}
In this section the degree \(H\) is the only quantity allowed to move.
The exponent, the weights and the centers are fixed. Let \(\nu>2\)
and let \(p_i,q_i\in\Z\) satisfy
\begin{equation}\label{eq:approximants}
 |\pi-p_i/q_i|\le q_i^{-\nu},\qquad q_i\ge2,\quad p_i\ne0,
 \qquad w_i=\lceil\log q_i\rceil,\quad r_i=2\ii p_i/q_i.
\end{equation}
Assume the weights satisfy Theorem~\ref{thm:interpolation}.
Use the centers \(c_{ji}=jr_i\), the multiples of the approximate periods
\(r_i\). They are coordinatewise distinct
because \(r_i\ne0\). Write
\[
 w_* =\min_{i>0}w_i,\qquad \omega=2\pi\ii,
 \qquad \Lambda=4\log2,\qquad c=\tfrac14\log2.
\]

The interpolation chart contains \(\log(1+t)\). For the determinant we
need polynomial entries, so choose \(F_0>2/\theta\), and replace the
logarithms by the polynomials
\begin{equation}\label{eq:truncation}
 T_i=\left\lceil F_0w_i/v_0\right\rceil,\qquad
 G_i(t)=\sum_{1\le k<T_i}\frac{(-1)^{k+1}}k t^k.
\end{equation}
An empty sum is zero. The first omitted term already has weight
\(v_0T_i\ge F_0w_i>w_i/\theta\). Adding this tail to \(u_i\), or
subtracting it, preserves each weighted cutoff ideal. The substitution
therefore induces an invertible change of coordinates on the jet space
and preserves the rank of the interpolation map.

Now read the interpolation map as a matrix. Its columns are the source
monomials \(P=Y^hX^\alpha\) of \(W\)-degree at most \(H\). Its rows are
the local coefficients that interpolation controls, indexed by
\begin{equation}\label{eq:rows}
 \rho=(j,s,\beta),\qquad 0\le j<K,\qquad
 v_0s+w\beta/\theta<H.
\end{equation}
Here \(\beta\), the multi-index of the \(u\)-coefficient, is the
transverse index. The entry is
\begin{equation}\label{eq:entry}
 [t^su^\beta]P(1+t,(jr_i+G_i(t)+u_i)_{i=1}^m).
\end{equation}
Surjectivity gives full row rank for large \(H\), so there are at least
as many columns as rows. Fix a nonzero square minor \(\Delta_H\) using
every row for the comparison at that \(H\). Let \(M=M_H\) be its size, and put
\begin{equation}\label{eq:rowcounts}
 \bar b_H=\frac1{MH}\sum_{\rho=(j,s,\beta)}w\beta,
 \qquad
 M\sim\frac{K\theta^mH^{m+1}}{(m+1)!v_0\prod_iw_i},
 \qquad 0\le\bar b_H\le\theta.
\end{equation}
The asymptotic follows by rescaling the fixed simplex by \(H\).
Its boundary has measure zero, so strict inequalities on that boundary
do not affect the leading lattice volume. The average \(\bar b_H\) records
transverse weight already present in the rows. It offsets part of the
column cost, leaving \(1-\bar b_H\) in the arithmetic lower bound.

\subsection{What the two cases must accomplish}
The three error terms come from denominator clearing, translation
and holomorphic growth.
\begin{align}
 E_{\rm ar}&=\frac{\Lambda F_0m}{v_0}
       +\Lambda\sum_i\frac1{w_i}+\frac{\theta}{w_*},\label{eq:Ear}\\
 E_{\rm tr}&=\frac{\nu}{F_0}+\frac{\log2}{v_0}
       +\frac{\log4+\log(2K)+\nu}{w_*},\label{eq:Etr}\\
 E_{\rm hol}&=\frac{100K}{w_0}+\frac{\log2}{v_0}
       +\frac{\log(200K)}{w_*},\qquad
 E_{\rm an}=E_{\rm tr}+E_{\rm hol}.\label{eq:Ehol}
\end{align}
Let \(A\in(0,1)\) be the cutoff that separates low transverse weight
from high transverse weight, and let
\(\eta\in(0,1)\) be the fraction of low rows that puts a term into the
first case rather than the second. The following proposition states the
entire comparison. The next three lemmas supply its three estimates:
the arithmetic lower bound, the translation to exact periods, and the
collision saving within one transverse index.

\begin{proposition}[Two-case comparison]\label{prop:comparison}
Define
\begin{equation}\label{eq:gL}
 g=\nu\bigl(A(1-\eta)-\theta\bigr)-(1-\theta),\qquad
 L=\frac{\eta^2K\theta^m}{(m+1)v_0A^m}.
\end{equation}
The data \eqref{eq:approximants} and the interpolation hypotheses cannot
hold simultaneously if
\begin{equation}\label{eq:contradiction}
 g>0,\qquad E_{\rm ar}+E_{\rm an}<g,\qquad
 cL>1+E_{\rm ar}+E_{\rm an}.
\end{equation}
\end{proposition}

For each choice in the row expansion, count the fraction of indices
at or below the cutoff. The third inequality of \eqref{eq:contradiction}
controls the case where that fraction is at least \(\eta\). The second
controls the other case.

\subsection{The arithmetic lower bound}
\begin{lemma}\label{lem:arithmetic}
Every selected nonzero minor satisfies
\begin{equation}\label{eq:lower}
 \frac{\log|\Delta_H|}{MH}\ge -(1-\bar b_H)-E_{\rm ar}.
\end{equation}
\end{lemma}
\begin{proof}
Set \(L_i=\lcm(1,\ldots,T_i-1)\), with \(L_i=1\) if \(T_i=1\),
and \(D_H=\prod_iL_i^{\lfloor H/w_i\rfloor}\).
Multiply column \(\alpha\) by \(\prod_iq_i^{\alpha_i}\) and row
\((j,s,\beta)\) by \(\prod_iq_i^{-\beta_i}\). The scaled entry is
zero unless \(\alpha\ge\beta\); otherwise it equals
\[
 \binom\alpha\beta[t^s](1+t)^h
       \prod_i\bigl(q_i(jr_i+G_i(t))\bigr)^{\alpha_i-\beta_i}.
\]
Here \(q_ijr_i=2\ii jp_i\in\Z[\ii]\). Its denominator divides
\(\prod_iL_i^{\alpha_i-\beta_i}\), which divides \(D_H\).
Multiplying every scaled entry by \(D_H\) therefore gives a matrix
over \(\Z[\ii]\) with nonzero determinant. Such a determinant has
modulus at least one. Undoing the scalings gives
\begin{equation}\label{eq:clearing}
 \log|\Delta_H|\ge -M\log D_H
 -\sum_{\text{columns }\alpha}\sum_i\alpha_i\log q_i
 +\sum_{\text{rows }\beta}\sum_i\beta_i\log q_i.
\end{equation}
The column sum is at most \(MH\). Since
\(0\le w_i-\log q_i<1\), the row sum is at least
\(MH\bar b_H-MH\theta/w_*\).

For completeness, \(\log\lcm(1,\ldots,n)\le4n\log2\).
If a prime power lies in \((\ell,2\ell]\), its logarithmic prime
weight contributes to \(\log\binom{2\ell}{\ell}\); the other
valuation contributions are nonnegative. Thus the increase of the log
least common multiple across this interval is at most \(2\ell\log2\).
Sum over dyadic intervals to the first power of two at least \(n\).
Consequently
\[
 \frac{\log D_H}{H}\le\Lambda\sum_i\frac{T_i}{w_i}
 \le\frac{\Lambda F_0m}{v_0}+\Lambda\sum_i\frac1{w_i}.
\]
Substitute these estimates into \eqref{eq:clearing}.
\end{proof}

\subsection{Translate rows to exact periods}
For every selected column \(P\), define the entire functions
\begin{equation}\label{eq:functions}
 f_{a,P}(z)=[u^a]P(e^z,z+u_1,\ldots,z+u_m).
\end{equation}
For \(P=Y^hX^\alpha\), this is
\(\binom\alpha a e^{hz}z^{|\alpha|-|a|}\) if \(a\le\alpha\),
and zero otherwise. Only indices with \(wa\le H\) can contribute.

\begin{lemma}[Exact translation]\label{lem:translation}
Each row \((j,s,\beta)\) is a sum of at most
\[
 Q_H=(\lfloor H\rfloor+1)^{2m}(\lfloor H/v_0\rfloor+1)
\]
scalar multiples of rows
\[
 \bigl([t^\ell]f_{a,P}(j\omega+\log(1+t))\bigr)_P,
 \qquad a\ge\beta,\quad wa\le H,\quad 0\le\ell\le s.
\]
Every scalar \(\xi\) with index \(a\) satisfies
\begin{equation}\label{eq:scalarbound}
 |\xi|\le\exp\{-\nu w(a-\beta)+HE_{\rm tr}\}.
\end{equation}
\end{lemma}
\begin{proof}
Put \(\epsilon_i=j(r_i-\omega)\) and
\(\tau_i=G_i-\log(1+t)\). For \(z=j\omega+\log(1+t)\), we have
\(e^z=1+t\) and \(z+\epsilon_i+\tau_i=jr_i+G_i\).
Expanding first in \(u\), then in \(\epsilon\) and \(\tau\), gives
the following identity of row vectors
\begin{align}
 \mathrm{row}_{j,s,\beta}
 ={}&\sum_{\substack{a\ge\beta\\wa\le H}}\binom a\beta
 \sum_{0\le d\le a-\beta}\binom{a-\beta}d\epsilon^{a-\beta-d}
 \sum_{k=0}^s[t^k]\prod_i\tau_i^{d_i}\notag\\[-2pt]
 &\hspace{23mm}\cdot
 \bigl([t^{s-k}]f_{a,P}(j\omega+\log(1+t))\bigr)_P.
 \label{eq:rowidentity}
\end{align}
In particular, the scalar factors do not depend on \(P\).

The tail product has order at least \(\sum_i d_iT_i\). A nonzero
coefficient in \eqref{eq:rowidentity} therefore satisfies
\begin{equation}\label{eq:tailbudget}
 wd\le\frac{v_0}{F_0}\sum_i d_iT_i
 \le\frac{v_0s}{F_0}<\frac H{F_0}.
\end{equation}
On \(|t|=1/2\), each \(|\tau_i(t)|\) is at most
\(\sum_{k\ge T_i}2^{-k}/k\le1\). Its product coefficient is thus
at most \(2^k\le2^s\) by Cauchy's estimate. Also
\[
 |\epsilon_i|\le2Kq_i^{-\nu}
 \le2Ke^\nu e^{-\nu w_i},\qquad
 \binom a\beta\binom{a-\beta}d\le4^{|a|},\qquad |a|\le H/w_*.
\]
Together with \eqref{eq:tailbudget}, these give
\[
 \log|\xi|\le-\nu w(a-\beta)+\frac{\nu H}{F_0}
 +\frac{H\log2}{v_0}
 +\frac{H(\log4+\log(2K)+\nu)}{w_*}
\]
for nonzero scalars; zero scalars satisfy the bound as well.
Finally \(w_i\ge1\), so each coordinate of \(a\) and \(d\) is at
most \(H\), and \(k\le H/v_0\). This gives the asserted term count.
\end{proof}

The logarithm tail can supply part of a translated coefficient
without a corresponding power of the approximation error. The term
\(\nu/F_0\) in the estimate accounts for that loss.

\subsection{The quadratic saving within one transverse type}
\begin{lemma}[Collision estimate]\label{lem:collision}
Choose one test row from each expansion in Lemma~\ref{lem:translation}
and remove its scalar. Let \(T\) be the resulting square matrix, and
let \(n_a\) count the rows with transverse index \(a\). There is a
quantity \(\rho_H\to0\), uniform in these choices and in the selected
columns, such that
\begin{equation}\label{eq:collision}
 |\det T|\le
 \exp\left\{-c\sum_a n_a^2+MH(E_{\rm hol}+\rho_H)\right\}.
\end{equation}
\end{lemma}
\begin{proof}
Set \(R=100K\). On \(|z|\le R\), formula \eqref{eq:functions}
and the monomial budgets give
\[
 |f_{a,P}(z)|\le
 \exp\left\{H\left(\frac R{w_0}+\frac{\log(2R)}{w_*}\right)\right\}
 =:U_H.
\]
Write \(f_{a,P}(z)=\sum_{d\ge0}c_{a,d,P}(z/R)^d\).
Cauchy's estimate gives \(|c_{a,d,P}|\le U_H\).
For \(|t|=1/2\),
\[
 |j\omega+\log(1+t)|\le2\pi(K-1)+\log2<R/2.
\]
Thus a test of order \(\ell\) applied to \((z/R)^d\) has modulus
at most \(2^\ell2^{-d}\).

Expand the determinant in these Taylor series. The expansion converges
absolutely: every coefficient determinant is bounded by
\(M!U_H^M\), and the product geometric series over its
finitely many rows is summable. Equal pairs \((a,d)\) give equal
coefficient vectors \((c_{a,d,P})_P\), so any term with a repeated
pair vanishes. For a nonzero term,
\[
 \sum_{\rm rows}d\ge\sum_a\binom{n_a}{2}.
\]
Use half the geometric decay to enforce this lower bound and sum the
other half freely. Since \(\sum\ell\le MH/v_0\),
\[
 |\det T|\le M!U_H^M2^{MH/v_0}
 2^{-\frac12\sum_a\binom{n_a}{2}}(1-2^{-1/2})^{-M}.
\]
Now \(\sum_an_a=M\), and \(\log M!\le M\log M\). Hence
\eqref{eq:collision} follows with
\[
 \rho_H=\frac{\log M+\frac14\log2-\log(1-2^{-1/2})}{H}.
\]
Equation~\eqref{eq:rowcounts} gives \(\rho_H\to0\).
\end{proof}

For example, three rows with the same transverse index cannot all
survive using degrees \(0,0,1\): the two degree-zero coefficient
rows are equal. Any surviving assignment to these rows has total
Taylor degree at least \(0+1+2=3\).
The general value \(n(n-1)/2\) is the source of the quadratic saving.
The argument requires repetition only within a transverse type,
not equality of all the functions appearing in the matrix.

\subsection{Finish the two-case comparison}
\begin{proof}[Proof of Proposition~\ref{prop:comparison}]
Let
\[
 D_A(H)=\#\{a\in\N^m:wa\le AH\},\qquad
 L_H=\frac{\eta^2M}{H D_A(H)}.
\]
The simplex count in dimension \(m\) gives
\(D_A(H)\sim(AH)^m/(m!\prod_iw_i)\); by \eqref{eq:rowcounts},
\(L_H\to L\). The factorial ratio is precisely \(1/(m+1)\).

Expand \(\Delta_H\) using \eqref{eq:rowidentity}. The product of
the scalars in any one choice is bounded by
\[
 \exp\left\{-\nu\sum_\rho wa_\rho+\nu MH\bar b_H+MH E_{\rm tr}\right\}.
\]
There are at most \(Q_H^M\) choices. Apply
Lemma~\ref{lem:collision} to each.

\emph{Case 1.} At least \(\eta M\) indices satisfy \(wa\le AH\).
There are at most \(D_A(H)\) such indices, so Cauchy--Schwarz gives
\[
 \sum_a n_a^2\ge\frac{(\eta M)^2}{D_A(H)}=MH L_H.
\]
Since \(a_\rho\ge\beta_\rho\), the scalar main exponent is
nonpositive. Discard it. After division by \(MH\), the logarithmic
upper bound is \(-cL_H+E_{\rm an}+\rho_H\).

\emph{Case 2.} Fewer than \(\eta M\) indices satisfy \(wa\le AH\).
Then \(\sum_\rho wa_\rho\ge MH A(1-\eta)\). Discard the
nonpositive collision saving. The upper bound is now
\(-\nu(A(1-\eta)-\bar b_H)+E_{\rm an}+\rho_H\).

Take absolute values and sum all terms. We obtain
\begin{equation}\label{eq:upper}
 \frac{\log|\Delta_H|}{MH}
 \le E_{\rm an}+\varepsilon_H+
 \max\{-cL_H,-\nu(A(1-\eta)-\bar b_H)\},
\end{equation}
where \(\varepsilon_H=\rho_H+\log Q_H/H\to0\).
The second case exceeds the arithmetic main cost by at least \(g\):
\[
 \nu(A(1-\eta)-\bar b_H)-(1-\bar b_H)
 =\nu A(1-\eta)-1-(\nu-1)\bar b_H\ge g.
\]
The first case has saving tending to \(cL\), and \(1-\bar b_H\le1\).
The strict inequalities \eqref{eq:contradiction} therefore make both
terms in \eqref{eq:upper} smaller than
\(-(1-\bar b_H)-E_{\rm ar}\) for sufficiently large \(H\).
This contradicts Lemma~\ref{lem:arithmetic} for the same minor.
\end{proof}

\section{Choose the parameters, then let the degree grow}\label{sec:parameters}
The threshold \(\nu>2\) enters when we choose the auxiliary parameters
for Proposition~\ref{prop:comparison}. The construction needs a rational
\(b\) with \(1/2<b<1-1/\nu\), an interval that is nonempty for positive
\(\nu\) exactly when \(\nu>2\). The lower endpoint leaves room for
\(A^2<\theta\), and the upper endpoint leaves room for the
approximation-error saving to exceed the arithmetic cost.

\begin{lemma}\label{lem:constants}
For every real \(\nu>2\), there are positive rational
\(\theta,A,B,C\) such that
\begin{equation}\label{eq:constants}
 0<\theta<A<B<1,\quad \nu(A-\theta)>1-\theta,
 \quad C>1,\quad B<1/C,\quad B<C\theta<1.
\end{equation}
One can then choose \(\eta\in(0,1)\) so that the number \(g\) in
\eqref{eq:gL} is positive.
\end{lemma}
\begin{proof}
Choose rational \(b\) with \(1/2<b<1-1/\nu\). For sufficiently
small positive rational \(\delta\), set
\(\theta=1-\delta\) and \(A=1-b\delta\). Then
\[
 \nu(A-\theta)-(1-\theta)=[\nu(1-b)-1]\delta>0,
 \qquad
 \theta-A^2=(2b-1)\delta-b^2\delta^2>0.
\]
Thus \((A/\theta,1/A)\) is nonempty. Choose rational \(C\) in that
interval and rational \(B\) between \(A\) and
\(\min(1/C,C\theta)\). We have \(C>1\) and
\(C\theta<\theta/A<1\), giving \eqref{eq:constants}.
Finally choose \(\eta>0\) small enough to preserve the strict gap.
\end{proof}

Lemma~\ref{lem:constants} enters the proof through the three ratios
\begin{equation}\label{eq:ratios}
 CB<1,\qquad C\theta/B>1,\qquad B/A>1.
\end{equation}
As the dimension \(m\) grows, \(CB<1\) makes the holomorphic-growth
cost \(K/w_0\) shrink, \(C\theta/B>1\) makes the time weight \(v_0\)
grow, and \(B/A>1\) makes the collision saving \(L\) grow.

\begin{proof}[Proof of Theorem~\ref{thm:main}, using Theorem~\ref{thm:interpolation}]
Fix \(\nu>2\), and suppose there are arbitrarily large positive
integers \(q\) with some \(p\in\Z\) satisfying
\(|\pi-p/q|\le q^{-\nu}\). Fix the constants from
Lemma~\ref{lem:constants}, and put \(\varepsilon_0=\min(g/2,1/2)\).

\emph{First choose the truncation budget.} Choose \(F_0>2/\theta\)
so large that \(\nu/F_0<\varepsilon_0/3\).

\emph{Next choose one finite dimension.} For a tentative integer \(m\),
put
\begin{equation}\label{eq:dimension}
 K=\lfloor C^m\rfloor,\qquad w_0=B^{-m},\qquad
 v_0=2K\theta^m w_0.
\end{equation}
These positive rational data satisfy
\[
 K\theta^m\le(C\theta)^m<1,\qquad
 K(w_0/v_0)\theta^m=\tfrac12.
\]
As \(m\to\infty\), the ratios \eqref{eq:ratios} give
\[
 \frac K{w_0}\le(CB)^m\to0,\qquad
 v_0\sim2(C\theta/B)^m\to\infty,\qquad \frac m{v_0}\to0.
\]
The collision quantity has the exact expression
\begin{equation}\label{eq:collisiongrowth}
 L=\frac{\eta^2(B/A)^m}{2(m+1)}\longrightarrow\infty.
\end{equation}
Choose \(m\) large enough that
\begin{equation}\label{eq:dimensionmargin}
 \frac{\Lambda F_0m+2\log2}{v_0}+\frac{100K}{w_0}
 <\frac{\varepsilon_0}{3},\qquad cL>2,
\end{equation}
and fix this \(m\), hence \(K,w_0,v_0\).

\emph{Next choose the approximation denominators.} Select finitely
many of the assumed good fractions \(p_i/q_i\), successively, so that
\(w_i=\lceil\log q_i\rceil\) satisfy the thresholds in
Theorem~\ref{thm:interpolation} and also
\begin{equation}\label{eq:weightmargin}
 \Lambda\sum_i\frac1{w_i}
 +\frac{\theta+\log4+\log(2K)+\nu+\log(200K)}{w_*}
 <\frac{\varepsilon_0}{3}.
\end{equation}
To do this, first impose a common lower bound on every \(w_i\), using
\(\sum_i1/w_i\le m/w_*\), and then impose the successive geometric
thresholds. The good denominators are unbounded, so every finite lower
requirement can be met. They need not occur densely or in prescribed
intervals. The weight thresholds do not depend on the centers.
For large denominators \(p_i\ne0\), making \(jr_i\) distinct in
each coordinate; no independence between coordinates is needed.

The complete error is
\begin{align*}
 E_{\rm ar}+E_{\rm an}
 &=\frac\nu{F_0}
 +\frac{\Lambda F_0m+2\log2}{v_0}+\frac{100K}{w_0}\\
 &\quad+\Lambda\sum_i\frac1{w_i}
 +\frac{\theta+\log4+\log(2K)+\nu+\log(200K)}{w_*}
 <\varepsilon_0.
\end{align*}
In particular it is less than \(g\), and
\(cL>2>1+E_{\rm ar}+E_{\rm an}\).
Proposition~\ref{prop:comparison} gives a contradiction. Its degree
limit is taken only now, after the dimension, weights and centers have
all been fixed. A center-dependent lower threshold for \(H\) is
therefore harmless.

For each fixed \(\nu>2\), good denominators are bounded. This yields
the stated \(Q(\nu)\) uniformly over all integer numerators.
We allowed unreduced fractions and zero error. If \(\pi\) were
rational, its exact representations with arbitrarily large denominators
would contradict the result; hence \(\pi\) is irrational.

For the reverse exponent inequality, place the \(N+1\) fractional
parts \(0,\pi,\ldots,N\pi\) in \(N\) equal half-open intervals.
Pigeonhole gives \(p\in\Z\), \(1\le q\le N\), with
\[
 0<|\pi-p/q|\le\frac1{qN}\le\frac1{q^2}.
\]
The errors tend to zero as \(N\to\infty\). Reducing the fractions
preserves the inverse-square bound for the reduced denominator.
Those denominators must be unbounded: otherwise only finitely many
reduced fractions in a fixed bounded interval could occur, incompatible
with positive errors tending to zero. Every exponent below two
therefore occurs infinitely often with strict error inequality. This
gives \(\mu(\pi)\ge2\), completing the proof once the interpolation
input is established.
\end{proof}

\begin{remark}[What the geometry must supply]
The argument has used only eventual surjectivity for the original
monomial space, with center-independent successive weight thresholds.
It has not used coefficients or heights of a constructed interpolating
basis. Thus an existence proof with a center-dependent degree threshold
is sufficient. The next two sections establish exactly this input.
\end{remark}
\section{The geometric obstruction is a curve inequality}\label{sec:curve}
It remains to prove Theorem~\ref{thm:interpolation}. The geometric
obstruction takes the form of a curve inequality: an algebraic curve
cannot follow the logarithmic charts closely at the prescribed centers
without a correspondingly large weighted degree. We first establish a
local comparison of transverse multiplicities, and then use a finite
chain of zero loci of derivatives to turn a violating curve into a
component to which that comparison applies.

\subsection{A local comparison between two normal bases}
At a smooth point of a codimension-\(k\) subvariety \(Z\subset\C^n\),
a set of \(k\) ambient vector fields is a \emph{normal basis} if their
images form a basis of the quotient of the ambient tangent space by
\(T Z\). All normal-basis assertions below are made on a common dense
smooth open subset of an irreducible \(Z\). The next lemma turns
vanishing of low-cost derivatives in one frame into a product inequality
for the weights of the two frames.

\begin{lemma}[Weighted transverse multiplicity]\label{lem:multiplicity}
Give affine coordinates \(x_1,\ldots,x_n\) positive rational weights
\(\rho_1,\ldots,\rho_n\). Let polynomials \(f_1,\ldots,f_r\) have
weighted degree at most \(N>0\), and let \(Z\) be an irreducible
codimension-\(k\) component of their common zero locus on a Zariski
open set. Let \(D_1,\ldots,D_n\) be commuting holomorphic vector
fields forming a frame there, with positive costs \(\kappa_i\).
Fix \(\epsilon>0\) and suppose
\[
 D^\gamma f_\ell|_Z=0
 \quad\text{whenever }\sum_i\kappa_i\gamma_i<\epsilon N,
 \qquad 1\le\ell\le r.
\]
If the coordinate fields indexed by \(A\) and the frame fields
indexed by \(B\) are normal bases, where \(|A|=|B|=k\), then
\begin{equation}\label{eq:multiplicity}
 \prod_{a\in A}\rho_a
 \le 2k!\epsilon^{-k}\prod_{b\in B}\kappa_b.
\end{equation}
The constant is independent of \(N\), \(Z\), and the chosen point.
\end{lemma}
\begin{proof}
Choose a smooth point \(x\) of \(Z\) where both normal-basis
conditions hold and no other component of the common zero locus
passes. Flow from a neighborhood in \(Z\) along the commuting
\(B\)-fields. The differential of this flow map is invertible at
\((x,0)\), so it gives holomorphic coordinates \((z,b)\), with
\(Z=\{b=0\}\). The Taylor coefficients in \(b\) of \(f_\ell\)
are the corresponding iterated derivatives restricted to \(Z\),
divided by factorials. All coefficients of cost less than
\(\epsilon N\) vanish identically along \(Z\).

Let \(S\) be the coordinate \(k\)-plane through \(x\) obtained by
varying the coordinates in \(A\). Its tangent space is complementary
to \(T_xZ\). The differential of the transverse map \(b\) has
kernel \(T_xZ\); hence its restriction to \(T_xS\) is invertible.
Thus the \(b\)'s are also holomorphic coordinates on \(S\) near \(x\).
Upon restricting \(f_\ell\) to \(S\), substitution in the analytic
coefficient functions cannot lower the cost of a surviving factor
\(b^\alpha\). Therefore the restricted ideal \(I\) is contained
in the monomial ideal of \(b\)-cost at least \(\epsilon N\).
The slice meets the local zero set only at \(x\), so \(I\) has
finite colength. It follows that
\begin{equation}\label{eq:lengthlower}
 \len(\OO_{S,x}/I)\ge
 \#\left\{\alpha\in\N^k:
       \sum_{b\in B}\kappa_b\alpha_b<\epsilon N\right\}
 \ge\frac{(\epsilon N)^k}{k!\prod_{b\in B}\kappa_b}.
\end{equation}
For the last inequality, floor the coordinates of every point in the
positive simplex. Its containing unit cube is anchored at an admissible
lattice point; the union of these cubes covers the simplex. The
simplex volume gives the lower bound, with no large-\(N\) assumption.

Choose \(k\) constant linear combinations of the restricted equations
whose ideal \(J\) has an isolated zero at \(x\). Such choices exist:
at each step avoid the finitely many proper coefficient subspaces
whose combinations vanish on a positive-dimensional local component
left by the previous equations. No such component is contained in all
generators of the maximal-ideal-primary ideal \(I\). The local
dimension drops at every step. Since \(J\subset I\),
\(\len(\OO_{S,x}/J)\ge\len(\OO_{S,x}/I)\).

We bound the former length by weighted B\'ezout. Choose an integer
\(M_0\) with \(M_0\rho_a\) positive integers, and pull back the
equations on \(S\) under
\[
 x_a=\lambda_a+z_a^{M_0\rho_a}\qquad(a\in A),
 \qquad \lambda_a\ne x_a(x).
\]
There are \(\prod_{a\in A}(M_0\rho_a)\) preimages of \(x\), all
unramified. Each has the same finite local quotient length. A
polynomial of weighted degree at most \(N\) pulls back to a polynomial
of ordinary degree at most \(M_0N\). The affine B\'ezout bound for
isolated zeros therefore gives
\begin{equation}\label{eq:lengthupper}
 \len(\OO_{S,x}/J)\le\frac{N^k}{\prod_{a\in A}\rho_a}.
\end{equation}
Other positive-dimensional components away from these isolated points
do not affect the upper bound. We use the isolated-zero form of
\cite[Corollary VIII.3]{mondal}; the local multiplicity is the local
quotient length in the smooth affine ambient space
\cite[Section IV.4.2]{mondal}. Analytic and algebraic local lengths
agree here, as both equal the length after completion.
Combining \eqref{eq:lengthlower} and \eqref{eq:lengthupper} proves
even the bound without the factor 2; the stated weaker version is
convenient below.
\end{proof}

\subsection{Separated weights and contact with a logarithmic chart}
Fix the initial data in \eqref{eq:volume}. Choose positive rational
\(\sigma<1\) so small that
\begin{equation}\label{eq:sigma}
 \begin{split}
 (1+3\sigma)^{m+1}K(w_0/v_0)\theta^m&<1,\\
 (1+3\sigma)^mK\theta^m&<1,\\
 (1+\sigma)\theta&<1.
 \end{split}
\end{equation}
Set \(L_0=m+2\), \(v_i=w_i/\theta\) for \(i>0\), and
\[
 C_* =\max_{1\le k\le m}2k!(L_0/\sigma)^k.
\]
Choose the \(w_i\) successively so that, for all subsets
\(A,B\subset\{0,\ldots,m\}\) of equal cardinality between 1 and \(m\),
\begin{equation}\label{eq:separation}
 \prod_{a\in A}w_a>C_*\prod_{b\in B}v_b
\end{equation}
whenever the largest positive index in \(A\mathbin\triangle B\)
belongs to \(A\).

These are genuinely successive lower bounds. If that largest index is
\(i\), the ratio of the two products is \(w_i\) times a positive
expression in the preceding weights and the fixed zeroth data. Every
higher index occurs in both sets or in neither; a common higher index
contributes \(w_j/v_j=\theta\). There are finitely many subset pairs.
A sufficiently large rational \(w_i\) enforces all inequalities with
largest differing index \(i\), and later choices leave them unchanged.
No center location occurs in these requirements.

Let \(C\) be an irreducible algebraic curve whose generic point is
in our affine chart, and take its smooth complete model. Define
\begin{equation}\label{eq:curvedegree}
 \deg_W C=\sum_P\max\left(0,
 -\frac{\ord_PY}{w_0},-\frac{\ord_PX_1}{w_1},\ldots,
 -\frac{\ord_PX_m}{w_m}\right).
\end{equation}
An identically zero coordinate has order \(+\infty\) and contributes
no pole. At each branch \(P\) through \(a_j=(1,c_{j1},\ldots,c_{jm})\),
put
\begin{equation}\label{eq:contact}
 h_P=\min\left(\frac{\ord_P(Y-1)}{v_0},
 \left(\frac{\ord_P(X_i-c_{ji}-\log Y)}{v_i}\right)_{i=1}^m\right).
\end{equation}
The logarithm is its formal expansion at \(Y=1\). The entries are
positive and possibly infinite. Their minimum is finite on a
nonconstant branch: either \(Y-1\) has finite order, or \(Y=1\)
and some \(X_i\) is nonconstant.

\begin{proposition}[Curve inequality]\label{prop:curve}
For the separated weights above and coordinatewise distinct centers,
every such curve satisfies
\begin{equation}\label{eq:CI}
 \deg_W C\ge(1+\sigma)\sum_{P\text{ over the centers}}h_P.
\end{equation}
The sum includes every branch at every prescribed center.
\end{proposition}

\subsection{From a violating curve to a persistent component}
\begin{proof}
Suppose \eqref{eq:CI} fails. The curve meets a center, since otherwise
the right side is zero. We first show that the curve lies in \(Y=1\)
and then that it meets at most one center. A comparison of zero and
pole orders at the end of the proof contradicts the assumed violation.

For all sufficiently large integers \(N\),
there is a nonzero \(F_N\in\Pcal_W(N)\) whose local order at each
center is at least \((1+3\sigma)N\). Indeed, the leading counts of
coefficients and homogeneous vanishing conditions are respectively
\[
 \frac{N^{m+1}}{(m+1)!w_0\prod_iw_i},\qquad
 \frac{K(1+3\sigma)^{m+1}\theta^mN^{m+1}}
 {(m+1)!v_0\prod_iw_i}.
\]
The first inequality of \eqref{eq:sigma} makes the first larger.
Only the number of conditions matters; no independence is assumed.

On \(Y\ne0\), use the commuting frame
\[
 D_0=Y\partial_Y+\sum_{i=1}^m\partial_{X_i},\qquad
 D_i=\partial_{X_i}\quad(i>0).
\]
These fields preserve the polynomial \(W\)-degree bound. In each
logarithmic chart they are \((1+t)\partial_t\) and
\(\partial_{u_i}\). Thus a derivative of total cost
\(\sum_av_a\gamma_a\le\sigma N\) retains local order at least
\((1+2\sigma)N\). On a branch \(P\), its vanishing order is at
least \((1+2\sigma)Nh_P\): substitute the branch into each local
monomial and use \eqref{eq:contact}.

Every polynomial of \(W\)-degree at most \(N\), restricted to the
complete curve, has total pole degree at most \(N\deg_W C\).
This follows monomial by monomial from \eqref{eq:curvedegree}; addition
cannot increase a pole bound. Our violating inequality gives strictly
more zeros than poles unless the restricted derivative is zero.
Consequently
\begin{equation}\label{eq:derivativezero}
 D^\gamma F_N|_C=0\qquad
 \left(\sum_av_a\gamma_a\le\sigma N\right).
\end{equation}

Put \(\delta=\sigma N/L_0\). For \(0\le r\le L_0\), let \(V_r\)
be the common zero set in \(Y\ne0\) of derivatives of cost at most
\(r\delta\). Each is algebraic, since only finitely many derivatives
have bounded cost. They all contain the irreducible open curve
\(C^\circ=C\cap\{Y\ne0\}\).
Choose an irreducible component of \(V_{L_0}\) containing this curve,
then components containing the chosen one backwards through the chain.
Their dimensions lie between 1 and \(m\), since \(F_N\ne0\).
A proper inclusion of irreducible varieties lowers dimension, so among
the \(L_0+1=m+3\) levels two adjacent chosen components agree in a
component \(Z\).

For the equations defining the earlier level, every further derivative
of cost less than \(\delta\) vanishes on \(Z\), because its total
cost is less than the next budget. Lemma~\ref{lem:multiplicity}, with
\(\epsilon=\sigma/L_0\), gives
\begin{equation}\label{eq:allbases}
 \prod_{a\in A}w_a\le C_*\prod_{b\in B}v_b
\end{equation}
for every coordinate normal basis \(A\) and frame normal basis \(B\)
of this component \(Z\), which contains \(C^\circ\) and has codimension
between \(1\) and \(m\).

\subsection*{The differential identity forces \(Y=1\)}
If every \(\partial_{X_i}\), \(i>0\), is tangent to \(Z\), they
span an \(m\)-dimensional subspace of its tangent space. Properness
forces equality, and hence \(dY|_Z=0\).
Otherwise let \(i>0\) be the greatest index with
\(\partial_{X_i}\) not tangent. Its nonzero normal image extends to
a coordinate normal basis \(A\) containing \(i\). All higher positive
directions are tangent, so they occur in no normal basis of either
frame. If a frame normal basis \(B\) omitted \(i\), the largest
positive index in \(A\mathbin\triangle B\) would be \(i\), contradicting
\eqref{eq:separation} and \eqref{eq:allbases}. Thus every frame normal
basis contains \(D_i\).

The other frame fields fail to span the normal quotient. Their ambient
span is the hyperplane
\[
 Q=\ker(dX_i-dY/Y).
\]
For a hyperplane \(Q\) and tangent space \(T\), its image in the
quotient by \(T\) fails to span exactly when \(T\subset Q\): otherwise
\(Q+T\) is the whole space. We obtain the identity of rational
differentials on \(Z\)
\begin{equation}\label{eq:logidentity}
 dX_i=dY/Y.
\end{equation}
This also forces \(Y\) to be constant. To see it without assuming
\(\dim Z=1\), suppose \(Y\) is nonconstant and choose a smooth point
where \(dY\ne0\) and \eqref{eq:logidentity} holds. Intersect with
\(\dim Z-1\) affine linear hyperplanes whose common tangent line
is not killed by \(dY\). The component through the resulting smooth
curve germ is an algebraic curve on which \(Y\) is nonconstant.
On its smooth complete model \eqref{eq:logidentity} is an equality
of rational differentials. A nonconstant meromorphic \(Y\) has a
zero or pole. Writing \(Y=s^nu\) there, with \(u\) a unit, gives
residue \(n\ne0\) for \(dY/Y\). An exact meromorphic differential
\(dX_i\) has zero residue, as is immediate from a Laurent expansion.
This is a contradiction.

In the first tangent-space case, vanishing \(dY\) also implies that
\(Y\) is constant, since the ground field has characteristic zero.
Thus in all cases \(Y\) is constant on \(Z\), and hence on \(C\).
The curve meets a prescribed center where \(Y=1\), so \(C\subset\{Y=1\}\).

\subsection*{The violating curve meets at most one center}
In the fiber \(Y=1\), repeat the coefficient count in \(m\) variables,
using the second inequality of \eqref{eq:sigma}. It gives a nonzero
polynomial \(G_N(X)\) of \(w\)-degree at most \(N\) and local order
at least \((1+3\sigma)N\) at all centers. Ordinary partial derivatives
of cost at most \(\sigma N\) again vanish on \(C\), by the same
zero-pole comparison.

If \(m=1\), this already contradicts nonzero \(G_N\): a curve in
the affine line is Zariski dense there. For \(m\ge2\), take the
same derivative-budget chain in the fiber. Its components containing
\(C\) have dimensions between 1 and \(m-1\). An adjacent repeated
component \(Z'\) therefore exists, and the local lemma yields
\eqref{eq:allbases}, now with basis sets in \(\{1,\ldots,m\}\).

Coordinate and frame normal bases now belong to the same family.
If there were two distinct basis sets, order them so that their
largest differing index belongs to \(A\); separation contradicts
\eqref{eq:allbases}. The basis set is therefore unique. It is nonempty.
Choose one of its indices \(i\). The remaining coordinate vectors
do not span the normal quotient, so the same hyperplane argument,
now with \(Q=\ker dX_i\), shows \(dX_i|_{Z'}=0\). Hence \(X_i\)
is constant on \(Z'\).

Coordinatewise distinctness now has its precise use: \(Z'\) contains
at most one prescribed center. Every prescribed center on \(C\) is
in \(C^\circ\subset Z'\), so \(C\) meets at most one, say \(a_j\).
Choose \(X_l\) nonconstant on \(C\). The rational function
\(X_l-c_{jl}\) has total pole degree at most \(w_l\deg_W C\),
and its zero order at each branch over \(a_j\) is at least
\((w_l/\theta)h_P\). Equality of zero and pole degrees implies
\[
 \deg_W C\ge\theta^{-1}\sum_Ph_P
 >(1+\sigma)\sum_Ph_P,
\]
using the last inequality of \eqref{eq:sigma}. This contradicts the
assumed violation and proves \eqref{eq:CI}. For \(K=1\) the
coordinatewise-distinctness condition is vacuous, and this argument
still applies.
\end{proof}

\section{From the curve inequality to interpolation}\label{sec:bridge}
The curve inequality bounds weighted source degree from below by a
strict multiple of total branch contact. On the ordinary blow-up of the
weighted ideals at the centers, this inequality makes the corresponding
combination of pullback degree and exceptional divisor nef, leaving room
to form an ample class that still weights contact strictly more heavily
than source degree. Global generation of a positive power yields
polynomials of bounded source degree and strictly greater local order.
Weighted arcs show that their initial forms have only the origin as a
common zero at each center. Each factor in homogeneous reduction gains
more local weight than it uses in source degree. Fixed Chinese remainder
lifts handle the low-weight remainders at bounded additional degree.
The resulting lifts are triangular across weights and centers, so
we can interpolate simultaneous jets with a strict linear degree
margin and a fixed additive cost.

\begin{samepage}
\begin{theorem}[Interpolation with a strict degree margin]\label{thm:bridge}
Let \(m\ge1\), let \(W,V\) be positive rational weight vectors, and
let \(a_j=(1,c_{j1},\ldots,c_{jm})\) be finitely many distinct centers.
Use the logarithmic charts and the definitions
\eqref{eq:curvedegree}--\eqref{eq:contact}. Suppose there is a rational
\(s>1\) such that every irreducible curve with generic point in the
affine chart satisfies
\begin{equation}\label{eq:generalCI}
 \deg_W C\ge s\sum_{P\text{ over the centers}}h_P.
\end{equation}
Then there are constants \(0<\kappa<1\) and \(B\ge0\), depending
on the fixed weights and centers, such that for every \(T>0\),
\begin{equation}\label{eq:margininterpolation}
 \Pcal_W(\kappa T+B)\longrightarrow\bigoplus_j\Jcal_V(T)
\end{equation}
is surjective. In particular, \eqref{eq:interpolation} is surjective
for every sufficiently large real \(H\).
\end{theorem}
\end{samepage}

We use four standard facts. The ordinary blow-up of an ideal
supported at finitely many points is universal for making the ideal
invertible. Its exceptional divisor is effective Cartier, and the
invertible sheaf of the negative exceptional divisor is relatively
ample. A relatively ample line bundle tensored with a sufficiently high
pullback power of an ample bundle is ample. On a projective complex
scheme, an ample rational divisor plus a nef rational divisor is ample.
Finally, an ample line bundle has a globally generated positive power.
These facts apply to our possibly singular projective schemes. References are
\cite[Sections 28.27, 29.38, 31.33]{stacks} and
\cite[Corollary 1.4.10]{lazarsfeld}.

\subsection{Replace the formal coordinates by polynomial coordinates}
To define the blow-up ideal in the affine polynomial ring, we truncate
the formal logarithm. The omitted terms must start at a weight above
that of the coordinate they modify. Choose a truncation \(G_i(t)\)
whose first omitted exponent \(T_i\) satisfies \(v_0T_i>v_i\).
At center \(j\), set
\[
 z_0=Y-1=t,\qquad z_i=X_i-c_{ji}-G_i(t).
\]
These are polynomial local coordinates, related to the logarithmic
ones by \(z_i=u_i+\log(1+t)-G_i(t)\). This substitution and its
inverse add only terms of weight greater than \(v_i\), so they
preserve every weighted jet cutoff. We use \(z\) for local weights and
initial forms. Source degrees are still measured in \(Y,X_1,\ldots,X_m\)
with weights \(W\).

The coordinate substitutions also preserve the branch contact \(h_P\).
If \(t\) is not
identically zero on a branch, the added tail has order divided by
\(v_i\) at least \(T_i\ord_Pt/v_i>\ord_Pt/v_0\). Terms beyond
that bound cannot change the minimum in \eqref{eq:contact}. If \(t\)
is identically zero, the tail vanishes. Hence in these coordinates
\begin{equation}\label{eq:algebraiccontact}
 h_P=\min_{0\le i\le m}\frac{\ord_Pz_i}{v_i}.
\end{equation}

\subsection{Use the curve bound to obtain global generators}
Choose \(R>0\) rational, sufficiently large, with every \(R/w_i\)
and \(R/v_i\) a positive integer. Embed \(\mathbb A^{m+1}\) in
projective space by all monomials of \(W\)-degree at most \(R\),
including 1, and let \(X\) be its projective closure. The monomials
include every original coordinate. The chart where the constant
coordinate is nonzero is therefore exactly the original affine space.
Let \(L=\OO_X(1)\).

At center \(j\), use the ideal
\[
 I_j=(z_0^{R/v_0},\ldots,z_m^{R/v_m}).
\]
These are finite-colength ideals supported at their respective centers.
They define a coherent ideal \(I\) on \(X\), equal to the unit ideal
off the centers. Let \(p:X'\to X\) be its ordinary blow-up, and write
\(I\OO_{X'}=\OO_{X'}(-E)\) for the exceptional divisor and
\(A=p^*L\) for the pullback of \(L\), in additive line-bundle
notation. The scheme \(X'\) is integral and projective, \(E\) is
effective Cartier, and \(-E\) is relatively ample.

For a curve meeting the affine chart, the monomial embedding has
\begin{equation}\label{eq:degreeidentity}
 \deg(L|_C)=R\deg_W C.
\end{equation}
Indeed, on its normalization the pole divisor needed for all embedding
coordinates has coefficient \(R\) times the maximum in
\eqref{eq:curvedegree}. A pure coordinate power attains each positive
maximum. When the maximum is zero, the constant monomial is nonvanishing.
After these poles are cleared the sections
have no common zero, which identifies their pullback line bundle and
its degree. Similarly, at a branch over a center,
\begin{equation}\label{eq:idealidentity}
 \ord_P I=\min_i(R/v_i)\ord_Pz_i=Rh_P.
\end{equation}

Recall that a divisor class is nef if it has nonnegative degree on every
integral curve. To show that \(A-sE\) is nef, consider first a
noncontracted integral curve in \(X'\) whose image \(C\) meets the
affine chart. The map to \(C\) is birational, so
\eqref{eq:degreeidentity} and \eqref{eq:idealidentity} give
\((A-sE)\)-degree
\[
 R\left(\deg_W C-s\sum_P h_P\right)\ge0
\]
by \eqref{eq:generalCI}. A curve mapping into
the boundary avoids \(E\) and has positive \(A\)-degree. A
contracted curve has zero \(A\)-degree and positive \((-E)\)-degree
by relative ampleness. These cases exhaust the integral curves,
so \(A-sE\) is nef.

Choose an integer \(a_0>1\) so large that \(a_0A-E\) is ample.
For any rational \(1<\tau<s\), the identity
\begin{equation}\label{eq:amplecombination}
 A-\tau E=
 \frac{a_0\tau-1}{a_0s-1}(A-sE)
 +\frac{s-\tau}{a_0s-1}(a_0A-E)
\end{equation}
expresses it as the sum of a positive multiple of a nef class and
a positive multiple of an ample class. It is ample. Clearing
denominators and taking a globally generated power yields integers
\(b>a>0\), with \(b/a=\tau\), and finitely many sections
\(s_1,\ldots,s_r\) generating \(\OO_{X'}(aA-bE)\).
For the last step, one may use the definition of ampleness by an
affine cover of nonvanishing sections: take a finite subcover and
raise its sections to a common degree. This produces global generation
without an ideal-power pushdown or a cohomological vanishing argument.

\subsection{The generators have a controlled polynomial degree}
Under the inclusion \(\OO(aA-bE)\subset\OO(aA)\), each section \(s_k\)
restricts on the affine chart away from the centers to a rational
function regular there. It extends to a polynomial \(P_k\).
Indeed, a nonconstant reduced denominator in
\(\C[Y,X_1,\ldots,X_m]\) would give a codimension-one pole locus;
it cannot be supported at finitely many points because \(m+1\ge2\).
On \(p^{-1}(\mathbb A^{m+1})\), this polynomial pullback agrees
with the section in the affine trivialization, by equality on a
dense open set and integrality.

We now prove the source degree bound
\begin{equation}\label{eq:generatorbudget}
 \deg_WP_k\le aR.
\end{equation}
Our earlier note~\cite{casperdegree} extends OpenAI's eventual degree
bound for sections on the monomial compactification to every bundle power
by taking tensor powers. Here we obtain \eqref{eq:generatorbudget}
directly on the blow-up by testing on a monomial curve. Choose a
positive integer \(q\) clearing all weights and take
\[
 (Y,X_1,\ldots,X_m)=(c_0h^{qw_0},c_1h^{qw_1},\ldots,c_mh^{qw_m}),
\]
with generic nonzero \(c_i\). The embedding pulls \(L\) back to
\(\OO_{\mathbb P^1}(qR)\): the constant coordinate prevents common
zeros at finite points, and a pure power of degree \(qR\) prevents
a common zero at infinity. The pullback of \(I\) is a nonzero ideal
on the smooth curve, hence invertible; the curve lifts to \(X'\).
The constant-coordinate section trivializes the pullback of \(L\) on the
finite \(h\)-line. Its \(a\)-th power vanishes to order
\(aqR\) at infinity, so the pulled-back section of \(aA\), expressed
in this trivialization, has pole order at most \(aqR\).
For generic \(c\), the highest weighted part of \(P_k\) does not
cancel, so its pole order is \(q\deg_WP_k\). This proves
\eqref{eq:generatorbudget}.

Set
\[
 d=aR,\qquad \lambda=bR,\qquad \kappa=d/\lambda=a/b<1.
\]
Each reduction factor has source degree at most \(d=\kappa\lambda\).
Every \(P_k\) has local order at least \(\lambda\) at every center.
Let \(f_{kj}\) be the weight-\(\lambda\) part of \(P_k\) at center \(j\).
At each center the common zero set of these parts is
\begin{equation}\label{eq:initialforms}
 \{z\in\C^{m+1}:f_{1j}(z)=\cdots=f_{rj}(z)=0\}=\{0\}.
\end{equation}
Weighted arcs give both properties. Clear the local
weights by \(q\), and for any nonzero \(c\in\C^{m+1}\) take
\(z_i=c_i h^{qv_i}\). Its pullback ideal is \((h^{qR})\), so it
lifts to \(X'\). Each section, viewed as a polynomial, vanishes
along this arc to order at least \(qbR\). A lowest-weight term
below \(bR\) would be nonzero on a generic direction \(c\),
contradicting this bound. At the special point of the lifted arc,
global generation makes one generating section a unit after the
factor \(h^{qbR}\) is removed. Therefore one \(f_{kj}(c)\) is
nonzero. This proves \eqref{eq:initialforms}, including directions
with some zero coordinates.

\subsection{Finish by finite polynomial reduction}
By \eqref{eq:initialforms} and the Nullstellensatz, the homogeneous ideal
\((f_{1j},\ldots,f_{rj})\subset\C[z_0,\ldots,z_m]\) has finite
colength. Choose a common number \(M_0>0\) so that, at every center,
each homogeneous polynomial of \(V\)-weight at least \(M_0\)
belongs to that ideal. Such a cutoff exists because a finite-dimensional
positively graded quotient has only finitely many nonzero weights.

Repeated homogeneous reduction expresses a monomial \(z^\alpha\)
of weight \(\omega\) as a finite linear combination of terms
\begin{equation}\label{eq:gradedreduction}
 f_{k_1j}\cdots f_{k_qj}z^\beta,
 \qquad \wt_V(\beta)<M_0,\qquad
 q\lambda+\wt_V(\beta)=\omega.
\end{equation}
At a reduction step, membership in the homogeneous ideal permits
homogeneous coefficients of weight \(\omega-\lambda\).
Each further step lowers that weight by \(\lambda>0\), so the
process terminates. For a monomial already below \(M_0\), take
\(q=0\).

There are only finitely many pairs \((j,\beta)\) with
\(\wt_V(\beta)<M_0\). The Chinese remainder theorem supplies
polynomials \(B_{j\beta}\) whose jets modulo weight at least
\(M_0\) equal \(z^\beta\) at center \(j\) and zero at the other
centers. To justify this in ordinary polynomial algebra, sufficiently
high ordinary maximal-ideal powers are contained in the weighted
cutoff ideals. The latter have distinct maximal radicals, hence are
pairwise comaximal. Let \(B\ge0\) bound the \(W\)-degrees of all
the finitely many \(B_{j\beta}\).

In \eqref{eq:gradedreduction}, replace each \(f_{kj}\) by \(P_k\)
and each \(z^\beta\) by \(B_{j\beta}\). Apply the same linear
combination to obtain a polynomial \(Q_{j\alpha}\). At center \(j\),
its part of weight \(\omega\) is exactly \(z^\alpha\), with no
lower terms. All other terms there have greater weight. At any other
center its order is greater than \(\omega\), since
\(q\lambda+M_0>\omega\). Each contributing product has source degree
bounded by
\[
 \deg_W(P_{k_1}\cdots P_{k_q}B_{j\beta})\le qd+B\le\kappa\omega+B.
\]
Therefore \(Q_{j\alpha}\) has that degree bound as well.

For a fixed \(T>0\), order the finitely many target monomials of
weight less than \(T\) by increasing weight, grouping equal weights
across all centers. The coefficient matrix of their
\(Q_{j\alpha}\)'s is triangular with identity diagonal blocks.
Every such polynomial belongs to \(\Pcal_W(\kappa T+B)\), so
inverting this finite matrix proves \eqref{eq:margininterpolation}.
Undoing the filtered coordinate change returns the logarithmic charts.
Taking \(T=H\) and \(H\ge B/(1-\kappa)\), with \(H>0\), proves
the final assertion of Theorem~\ref{thm:bridge}.

\begin{proof}[Completion of Theorem~\ref{thm:interpolation}]
For the initial data in \eqref{eq:volume}, choose the successive
weight thresholds in Section~\ref{sec:curve}. Proposition~\ref{prop:curve}
provides \eqref{eq:generalCI} with \(s=1+\sigma>1\), for every
coordinatewise distinct tuple of centers. Theorem~\ref{thm:bridge}
then supplies eventual interpolation. Its constants \(\kappa,B\)
may depend on that tuple, exactly as the theorem permits. All
successive weight thresholds were fixed independently of the tuple.
\end{proof}

\newpage
\section{Relation to the original proof and Lean formalization}
The cohomological argument can also yield a strict degree margin. If we
rescale the source weights by a rational factor \(u\) with \(1<u<s\),
the curve inequality remains strict with factor \(s/u\). Interpolation
at an admissible cutoff \(H\) then uses source degree at most \(H/u\)
in the original weights. Rounding up to a sufficiently large admissible
cutoff adds a bounded term. See \cite[Section~2.3]{openai} for the
argument with ordinary ideal powers and Serre vanishing, and
\cite[Section~5.4]{chudnovsky} for the version with moving
multiplicative centers.

The attached Lean companion \texttt{PiTwoCaseBridge}\,
\attachfile[icon=Paperclip,description={Lean companion and verification record},mimetype=application/zip]{PiTwoCaseBridge-2026-11-26.zip}
builds on OpenAI's formalization. Its declaration \texttt{theorem\_5\_1}
proves Theorem~\ref{thm:bridge}, and
\texttt{pi\_irrationalityExponent\_eq\_two} uses it in the two-case
determinant argument to prove \(\mu(\pi)=2\). Stock Lean and Nanoda
accepted the same export of
\texttt{PiTwoCaseBridge.main}. Its eventual lower bound is nonstrict.
For \(2<\nu'<\nu\) and sufficiently large \(q\) it gives
\[
 |\pi-p/q|\ge q^{-\nu'}>q^{-\nu},
\]
as required in Theorem~\ref{thm:main}.

\subsection*{Authorship and AI use}
OpenAI Astra agents assisted substantially with mathematical development,
Lean formalization, manuscript preparation, and review.
\begin{thebibliography}{9}
\bibitem{openai}
OpenAI, \emph{The irrationality exponent of \(\pi\) is 2},
manuscript dated September 24, 2026.
\href{https://github.com/openai/math/tree/fd4aeeb2ee4fc729c18d98444fed42fd0529eeeb/preprints/The-irrationality-exponent-of-pi-is-2-September-24-2026}{Manuscript and source}.

\bibitem{chudnovsky}
G.~V.~Chudnovsky, \emph{Irrationality exponents of logarithmic numbers:
Values, quotients, and graded approximation}, 2026.
\href{https://arxiv.org/abs/2610.28079v1}{arXiv:2610.28079v1}.

\bibitem{laurent1991}
M.~Laurent, Sur quelques r\'esultats r\'ecents de transcendance,
\emph{Ast\'erisque} \textbf{198--199--200} (1991), 209--230.
\href{https://www.numdam.org/item/AST_1991__198-199-200__209_0/}{Original article}.

\bibitem{laurent2000}
M.~Laurent, Interpolation determinants of exponential polynomials,
\emph{Publicationes Mathematicae Debrecen} \textbf{56} (2000),
457--473. \href{https://doi.org/10.5486/PMD.2000.2310}{doi:10.5486/PMD.2000.2310}.

\bibitem{casperdegree}
R.~M.~Casper, \emph{The weighted degree bound for monomial compactifications
holds at every power}, October 8, 2026.
\href{https://github.com/Mattie/math/tree/4430312c0422ee667a0bf8a0d26c48e1bdda51b5/preprints/The-weighted-degree-bound-for-monomial-compactifications-holds-at-every-power-October-8-2026}{Manuscript and proof package}.

\bibitem{mondal}
P.~Mondal, \emph{How many zeroes? Counting the number of solutions
of systems of polynomials via geometry at infinity}, Draft III,
November 5, 2021.
\href{https://arxiv.org/abs/1806.05346v5}{arXiv:1806.05346v5}.
The local length convention is in Section IV.4.2; the isolated-zero
B\'ezout bound is Corollary VIII.3, printed p.~144.

\bibitem{lazarsfeld}
R.~Lazarsfeld, \emph{Positivity in Algebraic Geometry I},
Ergebnisse der Mathematik und ihrer Grenzgebiete, vol.~48,
Springer, 2004. Corollary 1.4.10, printed p.~53.
\href{https://www.math.stonybrook.edu/~roblaz/Reprints/Laz.PAG.Chapt1.pdf}{Author's Chapter 1 text}.

\bibitem{stacks}
The Stacks Project Authors, \emph{The Stacks Project}, 2026.
\href{https://stacks.math.columbia.edu/tag/01PR}{Section 28.27}
(ample invertible sheaves),
\href{https://stacks.math.columbia.edu/tag/01VG}{Section 29.38}
(relative ampleness),
\href{https://stacks.math.columbia.edu/tag/01OF}{Section 31.33}
(blowing up),
\href{https://stacks.math.columbia.edu/tag/0806}{Lemma 31.33.5}
(universal property), and
\href{https://stacks.math.columbia.edu/tag/0BXX}{Section 53.2}
(curves and function fields).
\end{thebibliography}
\end{document}
